% Tradução brasileira revista a partir do workpass espanhol congelado;
% a fonte permanece em source_es.
% SGA 3, Exposé VIII, tradução brasileira: §7 até o Corolário 7.4.
% Autoridade: Polo--Gille Exp8-8nov09.pdf, pp. locais 23--26,
% leitor completo pp. 639--642.

\setcounter{footnote}{31}
\SGARefTarget{sga3:viii:section:7}\section*{7. Apêndice: sobre os monomorfismos de pré-esquemas em grupos}
\addcontentsline{toc}{section}{7. Apêndice: sobre os monomorfismos de pré-esquemas em grupos}

O resultado demonstrado nesta \SGARefLink{sga3:viii:section:7}{seção}
não é necessário no restante do seminário, salvo em X,
\SGARefLink{sga3:x:theorem:8-8}{8.8}, e nos Exposés XV e XVI. Depende de
maneira essencial do teorema de existência de grupos quocientes do Exposé
VI\(_A\).\footnote{Nota do editor: cf. Exposé VI\(_A\),
\SGARefLink{sga3:viii:corollary:3:2}{Teoremas 3.2} e
\SGARefLink{thm:sga3-VIA-3.3.2}{3.3.2}.}

O \SGARefLink{sga3:viii:corollary:5:7}{Corolário 5.7} suscita a questão
de quando se pode afirmar que um monomorfismo \(u:G\to H\) de
\(S\)-grupos é uma imersão, ou mesmo uma imersão fechada.

Vimos em VI\(_B\), \SGARefLink{cor:sga3-VIB-1.4.2}{1.4.2}, que isso ocorre
quando \(S\) é o espectro de um corpo \(k\), desde que \(G\) seja de
tipo finito sobre \(k\) e \(H\) seja localmente de tipo finito sobre \(k\).
Deduz-se facilmente que o mesmo resultado continua válido quando se
supõe simplesmente que \(S\) seja artiniano.\footnote{Nota do editor: devem
ser levadas em conta as adições feitas no Exposé VI\(_B\).}

Por outro lado, é fácil dar exemplos de monomorfismos bijetivos que não são
imersões, por exemplo, sobre a reta afim sobre um corpo ou sobre o
espectro de um anel de valorização discreta. Tomemos
\[
  H=(\mathbf Z/2\mathbf Z)_S,\qquad G=G_1\times_SG_2,
\]
onde \(G_1\) é o subgrupo aberto de \(H\) complementar ao ponto
fechado \(x\ne0\) na fibra \(H_s\simeq\mathbf Z/2\mathbf Z\), sendo
\(s\) um ponto fechado fixo de \(S\), e \(G_2\) é o subesquema fechado
de \(H\), união do subesquema dado pela seção identidade e do
subesquema fechado reduzido definido por \(x\). Verifica-se imediatamente
que \(G_2\) é estável pela multiplicação em \(H\) e, portanto, é um
esquema em grupos. As imersões \(G_i\to H\), \(i=1,2\), definem
então um homomorfismo de \(S\)-grupos
\[
  G=G_1\times_SG_2\longrightarrow H
\]
que é manifestamente um monomorfismo bijetivo, e até mesmo uma imersão
local, mas não é uma imersão. De fato, \(G\) e \(H\) são reduzidos;
\(G\) tem três componentes irredutíveis disjuntas, enquanto \(H\)
tem apenas duas.

Observe-se também que \(G_1\to H\) é um exemplo de subgrupo aberto de
\(H\) que não é fechado, em contraste com a situação dos grupos
algébricos sobre um corpo. A teoria das degenerações de curvas
elípticas fornece outros exemplos desse fenômeno nos quais \(H\) é liso sobre
\(S\), de dimensão relativa um, e \(G\) tem fibras conexas.

Poder-se-ia esperar\begingroup
\renewcommand{\thefootnote}{(\(\ast\))}%
\footnote[1]{De fato, M. Raynaud construiu um contraexemplo com \(G\) liso
e de fibras conexas; cf. XVI,
\SGARefLink{sga3:xvi:subsection:1:1:item:c}{1.1(c)}. Sem a hipótese de
fibras conexas, pode-se tomar \(S=\operatorname{Spec}(\mathbf Z_2)\),
\(G=(\mathbf Z/2\mathbf Z)_S\) privado de seu ponto fechado não identidade, e
\(H=(\mu_2)_S\).}%
\endgroup
que, supondo \(G\) plano sobre \(S\), e, por exemplo, \(G\) e \(H\) de
apresentação finita sobre \(S\), todo monomorfismo \(u:G\to H\) de
\(S\)-grupos fosse automaticamente uma imersão. Provaremos um resultado
desse tipo sob hipóteses adicionais.

Primeiro, podemos supor que \(S\) seja afim e, pela hipótese de
apresentação finita sobre \(G\) e \(H\), que \(S\) seja noetheriano: reduzimo-nos
ao espectro de um anel de tipo finito sobre
\(\mathbf Z\).\footnote{Nota do editor: cf. EGA~IV\(_3\), §8 e Exposé
VI\(_B\), §10.} Então \(G\) e \(H\) são noetherianos. Dizer que
\(u:G\to H\) é uma imersão fechada (respectivamente, uma imersão)
equivale então a dizer que \(u\) é um monomorfismo, o que se verifica
por hipótese, e é próprio (respectivamente, próprio em todo ponto de
\(u(G)\)).\footnote{Nota do editor: cf. EGA~IV\(_3\), 8.11.5 e 15.7.1.}

O critério valuativo da propriedade de ser próprio mostra que basta verificar, após
toda mudança de base \(S'\to S\) em que \(S'\) seja o espectro de um
anel de valorização discreta, completo se desejado, que o morfismo
\(u':G'\to H'\) possui a correspondente propriedade de ser próprio. O caso da
propriedade local foi omitido em EGA~II, 7.3, e deveria aparecer como errata em
EGA~IV.\begingroup
\renewcommand{\thefootnote}{(\(\ast\))}%
\footnote[0]{Cf. EGA~IV\(_3\), 15.7.}%
\endgroup
Ficamos reduzidos, portanto, ao caso em que o próprio \(S\) é o espectro de
um anel de valorização discreta completo, com a ressalva de que toda
hipótese adicional imposta a \(G,H,u\) seja estável por mudança de base.

Sejam \(s\) e \(s'\), respectivamente, o ponto fechado e o ponto genérico
de \(S\). Os homomorfismos induzidos nas fibras,
\[
  u_s:G_s\longrightarrow H_s,
  \qquad
  u_{s'}:G_{s'}\longrightarrow H_{s'},
\]
são imersões fechadas porque são monomorfismos de esquemas em grupos de
tipo finito sobre corpos (VI\(_B\),
\SGARefLink{cor:sga3-VIB-1.4.2}{1.4.2}). Podemos, portanto, identificar
\(G_{s'}\) com um subesquema fechado de \(H_{s'}\). Utilizaremos o
\SGARefLink{sga3:viii:lemma:7:1}{resultado seguinte}.

\medskip
\noindent\SGARefTarget{sga3:viii:lemma:7:1}\textbf{Lema 7.1.}
\textit{Seja \(S\) o espectro de um anel de valorização discreta, com
ponto genérico \(s'\). Seja \(H\) um \(S\)-pré-esquema e seja \(L'\) um
subpré-esquema fechado da fibra genérica \(H_{s'}\), e portanto também
um subpré-esquema de \(H\). Existe seu fecho esquemático \(L\) em
\(H\), o menor subpré-esquema fechado de \(H\) que contém \(L'\), cf.
EGA~I, 9.5. É também o único subpré-esquema fechado de \(H\), plano sobre
\(S\), cuja fibra genérica é \(L'\). A formação de \(L\) é funtorial
no par \((H,L')\) e comuta com os produtos fibrados sobre \(S\). Em
particular, se \(H\) é um \(S\)-grupo e \(L'\) é um subgrupo de
\(H_{s'}\), então \(L\) é um subgrupo de \(H\).}

A demonstração é imediata e fica a cargo do leitor.\begingroup
\renewcommand{\thefootnote}{(\(\ast\ast\))}%
\footnote[0]{Cf. EGA~IV\(_2\), 2.8.}%
\endgroup
Apliquemos isso ao par \((H,G_{s'})\). O monomorfismo \(u:G\to H\) fatora-se como
\[
  G\xrightarrow{\,u'\,}L\longrightarrow H,
\]
onde \(L\) é um subgrupo de \(H\), um subpré-esquema fechado plano sobre
\(S\), e \(u':G\to L\) induz um isomorfismo nas fibras genéricas. Assim,
\(u:G\to H\) é uma imersão (respectivamente, uma imersão fechada) se,
e somente se, \(u':G\to L\) o é. Ficamos reduzidos ao caso em que \(H\) é
plano sobre \(S\) e \(u_{s'}:G_{s'}\to H_{s'}\) é um isomorfismo, novamente
com a exigência de que toda hipótese adicional continue válida
ao substituir \(H\) por um subpré-esquema em grupos fechado.

Como \(H\) é então o fecho esquemático de \(H_{s'}\), se \(u\)
é uma imersão, sua imagem \(u(G)\) é um subpré-esquema\footnote{Nota do
editor: suprimiu-se a palavra «fechado».} de \(H\) que contém
\(H_{s'}\). Seu fecho esquemático é, portanto, \(H\), de modo que
\(u(G)\) é um subpré-esquema aberto de \(H\). Consequentemente, devemos
provar que \(u\) é uma imersão aberta, ou um isomorfismo se a conclusão
desejada é que seja uma imersão fechada.

Como \(G\) e \(H\) são planos sobre \(S\), isso equivale a que os
morfismos induzidos nas fibras sejam imersões abertas,
respectivamente isomorfismos (cf. SGA~1, I,
\SGARefLink{sga3:viii:corollary:5:7}{5.7}). Isso já se verifica para \(u_{s'}\),
de modo que resta provar que \(u_s\) é uma imersão aberta,
respectivamente um isomorfismo.

Como \(G\) e \(H\) são planos sobre \(S\), as dimensões de suas fibras
são constantes (VI\(_B\),
\SGARefLink{sga3:viii:corollary:4:3}{4.3}). Suas fibras genéricas coincidem,
logo essas dimensões são iguais para \(G\) e \(H\). Assim,
\(u_s:G_s\to H_s\) é um monomorfismo de grupos algébricos da mesma
dimensão. Segue-se que \(G_s\) é aberto em \(H_s\), e é, como conjunto,
uma união de componentes conexas de \(H_s\). De fato, após passar a um
corpo-base algebricamente fechado e reduzir \(G\) e \(H\), ambos são
lisos sobre o corpo e a afirmação é imediata.

O complemento \(H_s-G_s\) é fechado em \(H_s\), e portanto em \(H\).
Seu complemento \(H'\) em \(H\) é aberto e estável pela lei de grupo.
Substituindo \(H\) por \(H'\), podemos reduzir-nos, com a ressalva habitual,
ao caso em que \(u\) é bijetivo, isto é, \(u_s:G_s\to H_s\) é
bijetivo. Em todos os casos ficamos reduzidos a provar que \(u\), ou,
equivalentemente, \(u_s\), é um isomorfismo, eventualmente sob essa
hipótese de bijetividade.

Suponhamos primeiro que \(u\) seja bijetivo. Se \(H_s\) é reduzido,
então \(u_s\) é um isomorfismo, porque \(G_s\) se identifica com um
subesquema fechado de \(H_s\) que tem o mesmo conjunto subjacente. Em
particular, se \(k=\kappa(s)\) é de característica zero, todo grupo
algébrico sobre \(k\) é reduzido pelo teorema de Cartier (cf.
VI\(_B\), 1.6.1; VII\(_B\), 3.3.1; ou EGA~IV\(_4\), 16.12.2 e 17.12.5).
Provamos:

\medskip
\noindent\SGARefTarget{sga3:viii:proposition:7:2}\textbf{Proposição 7.2.}
\textit{Seja \(u:G\to H\) um homomorfismo de pré-esquemas em grupos de
apresentação finita sobre \(S\). Se \(u\) é um monomorfismo, \(G\) é
plano sobre \(S\), e todo corpo residual de \(S\) tem característica
zero, então \(u\) é uma imersão.}

Suponhamos agora que \(\kappa(s)\) tenha característica \(p>0\), e
restrinjamo-nos ao caso em que \(G\) é comutativo. Sob as reduções
anteriores, \(H\) também é comutativo, porque é o fecho
esquemático de \(H_{s'}\simeq G_{s'}\). Para \(n\geq0\), ponhamos
\[
  S_n=\operatorname{Spec}(V/\mathfrak m^{n+1}),\qquad
  G_n=G\times_SS_n,\qquad H_n=H\times_SS_n,
\]
onde \(V\) é o anel de valorização que define \(S\) e \(\mathfrak m\)
seu ideal maximal. Para \(m>0\), denotemos por \({}_mG\) e \({}_mH\) os
núcleos da multiplicação por \(m\) em \(G\) e \(H\). Definamos de modo
análogo
\[
  {}_m(G_n)=({}_mG)_n,
\]
que abreviamos por \({}_mG_n\), e do mesmo modo para \(H\).

Por VI\(_A\), \SGARefLink{sga3:viii:corollary:3:2}{3.2}, existem os
quocientes
\[
  Q_n=H_n/G_n.
\]
Cada \(Q_n\) é um esquema em grupos comutativo sobre \(S_n\), plano sobre
\(S_n\), e \(H_n\to Q_n\) é fielmente plano com núcleo \(G_n\). A
formação do quociente comuta com a extensão da base,\footnote{Nota
do editor: isso deveria ser explicitado nos Exposés V e VI\(_A\).} de
modo que
\[
  Q_n\simeq Q_m\times_{S_m}S_n\qquad(m\geq n).
\]
Em particular, \(Q_n\times_{S_n}S_0=Q_0=H_0/G_0\). Como \(G_0\) e
\(H_0\) têm o mesmo conjunto subjacente, \(Q_0\) é, como conjunto, um
único ponto: é um grupo puramente infinitesimal.

Segue-se que cada \(Q_n\) é finito e plano sobre \(S_n\). Seja definido por
uma álgebra \(C_n\) sobre \(V_n=V/\mathfrak m^{n+1}\), livre finita como
\(V_n\)-módulo. Para \(m\geq n\),
\[
  C_n=C_m\otimes_{V_m}V_n,
\]
de maneira compatível também com os morfismos de comultiplicação. Portanto,
\[
  C=\varprojlim C_n
  \quad\text{é um módulo livre finito sobre}\quad
  V=\varprojlim V_n,
\]
e as comultiplicações dos \(C_n\) induzem uma sobre \(C\). Assim,
\(Q=\operatorname{Spec}(C)\) é um esquema em grupos finito e plano sobre
\(S\) que satisfaz
\[
  Q\times_SS_n=Q_n
\]
para todo \(n\).

\medskip
\noindent\SGARefTarget{sga3:viii:lemma:7:3}\textbf{Lema 7.3.}
\footnote{Nota do editor: eliminou-se aqui a hipótese de que \(Q\)
seja comutativo, e modificou-se em conformidade
\SGARefLink{sga3:viii:remark:7:3:1}{7.3.1}. Se \(Q\) não é comutativo, o
morfismo potência \(n\)-ésima não precisa ser um homomorfismo de
grupos.}
\textit{Seja \(K\) um corpo e seja \(Q\) um esquema em grupos finito sobre
\(K\), de grau \(n\). Então o morfismo potência \(n\)-ésima sobre
\(Q\) é nulo.}

Veja-se VII\(_A\), \SGARefLink{sga3:VIIA:8.5}{8.5}.

\medskip
\noindent\SGARefTarget{sga3:viii:remark:7:3:1}\textbf{Observação 7.3.1.}
O enunciado do Lema~\SGARefLink{sga3:viii:lemma:7:3}{7.3} continua tendo
sentido para um esquema em grupos finito localmente livre \(Q/S\) sobre um
pré-esquema de base arbitrário \(S\). Seria interessante encontrar uma demonstração
nessa generalidade.

O Lema~\SGARefLink{sga3:viii:lemma:7:3}{7.3} (isto é, VII\(_A\), 8.5)
mostra que o enunciado proposto é válido sempre que \(S\) seja
reduzido: basta aplicá-lo às fibras de \(Q\) nos pontos maximais de
\(S\), equivalentemente, nos pontos genéricos de suas componentes
irredutíveis.\footnote{Nota do editor: isso deveria ser acrescentado ao Exposé
VII\(_A\). Por outro lado, o enunciado é verdadeiro para todo \(S\) quando
\(Q\) é comutativo; cf. o teorema de Deligne em
\SGARefLink{sga3:viii:bib:to70}{[TO70]}, p.~4.}

Na situação anterior, em que \(S\) é o espectro de um anel de
valorização discreta e \(Q\) é comutativo, segue-se que
\(n\operatorname{id}_Q=0\). Aqui \(n\) é uma potência
\(p^{\nu_0}\) da característica residual. Portanto:

\medskip
\noindent\SGARefTarget{sga3:viii:corollary:7:4}\textbf{Corolário 7.4.}
\textit{Com \(Q\) e os \(Q_n\) como antes, suponhamos que seu grau comum
seja \(p^{\nu_0}\). Então}
\[
  p^{\nu_0}\operatorname{id}_Q=0,
  \qquad
  p^{\nu_0}\operatorname{id}_{Q_n}=0
  \quad\textit{para todo \(n\)}.
\]
